\section{The conditional moment induction}\label{sec:induction}
\label{ac:part:148}


We prove Theorem~\ref{thm:conditional} by induction on the number \(b\)
of coordinates, simultaneously for every allowed grid, path family, and
representation. First fix the constants in the order required by the
sparse estimate. Take \(\theta=a/16\) and \(K_0=10/e_0\) in
Lemma~\ref{lem:sparse}. Its exponent \(\rho\) is uniform for sufficiently
large \(R\), so we may next choose an integer
\[
 q\ge\max\{1,\lceil2/\rho\rceil\},\qquad 2q\rho\ge4,
\]
and only then choose a sufficiently large power of two \(R\). Every
largeness requirement below will hold uniformly for \(s\ge R\). In
particular, require
\[
 \frac1{\sqrt{\log R}}\le\frac{e_0}{2},
 \qquad
 \frac1{\log R}\le\frac{e_0}{4}.
\]
Since \(b\le\log s/\log R\) and \(C(\mathcal H)\le bh\), these choices give
\[
 c(s)\le2c_0,\qquad e(s)\ge e_0/2,\qquad
 e(s)h\log s-C(\mathcal H)\ge(e_0/4)h\log s.
\]
Further increases of \(R\) will absorb the displayed errors below.
After \(R\) is fixed, choose \(z_*>0\) small enough for
\eqref{ac:eq:3}, Lemma~\ref{lem:sparse}, and the base case.

Every conditional average is a contraction. Thus when \(F(\lambda)=0\),
its moment is at most one and the nonnegative allowance above proves the
theorem. This covers the empty diagram, \(M=1\), and all one-dimensional
types. A zero moment satisfies the bound by the convention
\(\log0=-\infty\).

For \(b=1\) and \(z=0\), conditioning on the prescribed assignments
leaves a uniform bijection of the remaining slots. Every nontrivial
average is therefore zero. Once \(R\) is fixed there are only finitely
many one-coordinate grids, path families, and representations, and their
conditional operators are continuous at zero. Decreasing \(z_*\) gives
the base case on a common interval.

Now assume \(b\ge2\) and the theorem for fewer coordinates. The contraction
bound gives \(T_\lambda^{\mathcal H}\le D_\lambda\), so the desired
inequality is automatic whenever
\begin{equation}
 (1+c(s))F(\lambda)\le e(s)h\log s-C(\mathcal H).
 \label{ac:eq:16}
\end{equation}
Suppose \(F(\lambda)<s^{1-a/4}\) and \eqref{ac:eq:16} fails. The positive
allowance above then gives
\(h\log s\le K_0F(\lambda)\), and \(M=s-h\ge s/2\) for large \(R\).
Put \(k=\min(M-\lambda_1,M-\lambda'_1)\). The dimension estimates from
Section~\ref{sec:prelim} give
\(c_{\rm dim}k\le F(\lambda)\le k\log s\). Hence \(h\le K_0k\), and
\[
 h+k\le s^{1-\theta}
\]
for large \(R\), because \(\theta=a/16<a/4\). Also \(k\ge1\), since
\(F(\lambda)>0\). Lemma~\ref{lem:sparse} applies. There are
\(D_\lambda\) singular values, so
\[
 \log T_\lambda^{\mathcal H}
 \le F(\lambda)-2q\rho k\log s
 \le-3F(\lambda)
 \le-c(s)F(\lambda)+e(s)h\log s-C(\mathcal H).
\]
This settles the sparse-dimension range.

We are left with \(b\ge2\) and the large-dimension range
\begin{equation}
 F(\lambda)\ge s^{1-a/4}.
 \label{ac:eq:17}
\end{equation}

\subsection{The hook estimate on the current grid}
\label{ac:part:199}

Removing the boxes outside a hook will record extra card trajectories.
We therefore first prove an estimate that is uniform over the resulting
path families. Set
\[
 p=\lfloor s^a\rfloor,\qquad
 \widehat c=c(s)+\frac{1}{20\sqrt{\log s}},\qquad
 \widehat e=e(s)-\frac{1}{20\sqrt{\log s}}.
\]
Under the induction hypothesis for fewer than \(b\) coordinates, we claim
that every family \(\mathcal H'\) of \(h'\) prescribed trajectories on
the current grid, with distinct slots at every layer and satisfying the
coordinate rule of Section~\ref{sec:paths}, and every \((p,p)\)-hook
diagram \(\gamma\vdash s-h'\), satisfy
\begin{equation}
 \log T_\gamma^{\mathcal H'}
 \le-\widehat c F(\gamma)+\widehat e h'\log s
       -C(\mathcal H')+s^{9/10}.
 \label{ac:eq:21}
\end{equation}
There is no lower bound on \(F(\gamma)\) in this claim. It retains the full
trajectory potential and improves both coefficients slightly. We will
derive it from two shorter coordinate sweeps, then use its dimension
reserve when the removed boxes are restored.

Split the coordinate list into consecutive parts of \(\lfloor b/2\rfloor\)
and \(\lceil b/2\rceil\) coordinates, with products \(s_1\) and \(s_2\).
Each coordinate contributes between \(\log R\) and \(2\log R\), so
\[
 \frac15\log s\le\log s_u\le\frac45\log s
 \qquad(u=1,2).
\]
In particular \(p\le s_u\). Regard the grid as \(s_2\) rows of length
\(s_1\). The first part of the sweep acts within rows and the second
within columns. At their junction, the middle positions of \(\mathcal H'\)
are removed, leaving a fixed board of \(s-h'\) free sites. Splitting each
prescribed trajectory there gives the path constraints in each row and
each column.

Use the signed tensor representation of Section~\ref{sec:density} on the
free sites at each layer. Let \(X\) be the conditional row average from
the input layer to the junction, and \(Y\) the conditional column average
from the junction to the output. The row and column randomizations are
conditionally independent under the fixed path constraints, so \(YX\) is
the full conditional average. At the junction let \(P_\gamma\) be the
full-group isotypic projection, and let \(P_\alpha^R,P_\beta^C\) be the
row and column subgroup projections. Here \(\alpha_i\) partitions the
number of free sites in row \(i\), and \(\beta_j\) partitions the number
of free sites in column \(j\); a row or column with no free sites has the
empty type. The sums below range over the local types occurring in the
signed tensor representation. Put
\[
 F_R=\sum_iF(\alpha_i),\qquad F_C=\sum_jF(\beta_j).
\]
The full projection commutes with both subgroup projections. Bijections
between the layers intertwine their full isotypic projections. Thus,
under compatible unitary identifications, the restriction of
\(Y P_\gamma X\) from the input \(\gamma\)-isotypic space to the output
\(\gamma\)-isotypic space is \(K_\gamma^{\mathcal H'}\otimes I\), where
the identity acts on the tensor multiplicity space. The hook type occurs
by Lemma~\ref{lem:density}, so this multiplicity is at least one.
Inserting the subgroup decompositions at the junction gives
\begin{equation}
\begin{aligned}
 (T_\gamma^{\mathcal H'})^{1/(2q)}
 &\le\|Y P_\gamma X\|_{2q}\\
 &\le\sum_{\alpha,\beta}
  \|Y P_\beta^C\|_{2q}\,
  \|P_\beta^C P_\gamma P_\alpha^R\|_{\rm op}\,
  \|P_\alpha^R X\|_{2q}.
\end{aligned}
\label{ac:eq:20}
\end{equation}
The first inequality uses the unnormalized trace on at least one copy of
\(\gamma\). The second uses the triangle inequality and the product norm
inequality, followed by \(\|A\|_{\rm op}\le\|A\|_{2q}\) for an outer
factor. We now bound the two child moments and the middle overlap.

Ordering the free sites by rows at both ends of the first part makes
\(X=\bigotimes_i X_i\); signed unitary reordering justifies this tensor
product. Write \(\mathcal H_i\) for the path constraints in row \(i\)
and \(\operatorname{mult}(\alpha_i)\) for the tensor multiplicity. Then
\[
 \|P_\alpha^R X\|_{2q}^{2q}
 =\prod_i \operatorname{mult}(\alpha_i)T_{\alpha_i}^{\mathcal H_i}.
\]
The column formula is identical. Lemma~\ref{lem:density} bounds each
multiplicity by \(\exp(O(p^2\log(s+1)))\). Apply the induction hypothesis
to every local constrained sweep and put
\[
 c_*=\min_{u=1,2}c(s_u),\qquad
 e_*=\max_{u=1,2}e(s_u).
\]
Each stage line belongs to exactly one child subproblem, so their
potentials add to \(C(\mathcal H')\). Each part has \(h'\) prescribed
trajectories, and their size weights add as
\(\log s_1+\log s_2=\log s\). Hence
\begin{equation}
\begin{aligned}
 \log\!\left(
 \|P_\alpha^R X\|_{2q}^{2q}\|Y P_\beta^C\|_{2q}^{2q}
 \right)
 \le{}&-c_*(F_R+F_C)+e_*h'\log s-C(\mathcal H')\\
 &+O((s_1+s_2)p^2\log(s+1)).
\end{aligned}
\label{ac:eq:18}
\end{equation}

The middle factor in \eqref{ac:eq:20} depends only on the fixed free-site
board. Its estimate does not use the trajectory probabilities.

\begin{lemma}[Overlap at the junction]\label{lem:overlap}
For the row and column type lists on the free junction sites and the full
hook type \(\gamma\) defined above,
\begin{equation}
 \log\|P_\beta^C P_\gamma P_\alpha^R\|_{\rm op}^2
 \le F_R+F_C-F(\gamma)+h'
       +O((s_1+s_2)p^2\log(s+1)).
 \label{ac:eq:19}
\end{equation}
Here a zero norm has logarithm \(-\infty\), and the constant is absolute.
\end{lemma}
\begin{proof}
Since \(P_\gamma\) commutes with both subgroup projections,
\[
 \|P_\beta^C P_\gamma P_\alpha^R\|_{\rm op}^2
 \le\|P_\beta^C P_\gamma P_\alpha^R\|_{\rm HS}^2
 =\Tr(P_\alpha^R P_\beta^C P_\gamma).
\]
The product \(P_\beta^C P_\gamma\) is a positive projection. We may
therefore dominate \(P_\alpha^R\) inside this trace by
Lemma~\ref{lem:density}. Apart from the factor
\(\exp(F_R+O(s_2p^2\log(s+1)))\), this leaves a probability mixture
of product densities whose one-site matrix in row \(i\) is an even
density \(r_i\).

For one such row family, average over all \(s_2\) rows, including empty
ones. For \(0<\varepsilon<1\), put
\[
 \bar r=\frac1{s_2}\sum_i r_i,\qquad
 \bar r_\varepsilon=(1-\varepsilon)\bar r+\varepsilon I/\dim V,
 \qquad A_\varepsilon=\bar r_\varepsilon^{-1/2}.
\]
The regularized reference density is even. With \(n=s-h'\), the positive
operator \(B_\varepsilon=\bar r_\varepsilon^{\otimes n}\) has trace one
and commutes with the full signed group action. On the
\(\gamma\)-isotypic space it has the form \(I_{D_\gamma}\otimes B_\gamma\).
Every eigenvalue there is therefore counted at least \(D_\gamma\) times
in its trace, and
\[
 B_\varepsilon P_\gamma\preceq D_\gamma^{-1}P_\gamma.
\]
It also commutes with \(P_\beta^C\). Compress this inequality by that
projection, use \(P_\beta^C P_\gamma\preceq P_\beta^C\) on the right, and
conjugate by \(B_\varepsilon^{-1/2}\). The result is
\[
 P_\beta^C P_\gamma
 \preceq D_\gamma^{-1}
 A_\varepsilon^{\otimes n}P_\beta^C A_\varepsilon^{\otimes n}.
\]

Now dominate the column projection by its density mixture. Apart from
the corresponding factor
\(\exp(F_C+O(s_1p^2\log(s+1)))\), each resulting mixture integrand is
the tensor-product trace
\[
 \prod_{(i,j)\text{ free}}
 \operatorname{tr}(A_\varepsilon r_i A_\varepsilon v_j),
\]
where \(v_j\) is the even density for column \(j\). Each factor is
nonnegative: it equals
\(\|r_i^{1/2}A_\varepsilon v_j^{1/2}\|_{\rm HS}^2\).
For every column, summing these factors over all rows gives
\[
 \sum_i\operatorname{tr}(A_\varepsilon r_i A_\varepsilon v_j)
 =s_2\operatorname{tr}(A_\varepsilon\bar r A_\varepsilon v_j)
 \le\frac{s_2}{1-\varepsilon}.
\]
If the column is missing \(l<s_2\) sites, AM--GM on its \(s_2-l\)
retained sites bounds their product by
\[
 (1-\varepsilon)^{-(s_2-l)}
 \left(\frac{s_2}{s_2-l}\right)^{s_2-l}
 \le(1-\varepsilon)^{-(s_2-l)}e^l.
\]
If \(l=s_2\), the empty column contributes the empty product one.
There are \(h'\) missing sites in total. Multiplication over columns,
averaging the uniformly bounded integrands, and then letting
\(\varepsilon\downarrow0\) prove \eqref{ac:eq:19}.
\end{proof}

We now combine the two estimates. Write \(W\) for the right side of
\eqref{ac:eq:19}, including its error. Terms with zero middle factor in
\eqref{ac:eq:20} contribute nothing. For a nonzero middle factor, its norm
is at most one as well as satisfying \eqref{ac:eq:19}. Since
\(c_*\le2c_0<1\le q\),
\[
 q\log\|P_\beta^C P_\gamma P_\alpha^R\|_{\rm op}^2
 \le c_*\min(0,W)\le c_*W.
\]
Adding this to \eqref{ac:eq:18} cancels \(F_R+F_C\). The number of lists
of local hook types contributes only
\(O_q((s_1+s_2)p^2\log(s+1))\) when the sum in \eqref{ac:eq:20} is
raised to \(2q\). We obtain
\[
 \log T_\gamma^{\mathcal H'}
 \le-c_*F(\gamma)+e_*h'\log s-C(\mathcal H')
       +c_*h'+O_q((s_1+s_2)p^2\log(s+1)).
\]

The size split gives the reserves in the claimed coefficients. With
\(L=\log s\),
\[
 c_*-c(s)=e(s)-e_*
 \ge\frac{\sqrt5/2-1}{\sqrt L}
 >\frac1{10\sqrt L}.
\]
The hats use only \(1/(20\sqrt L)\) of each reserve. The remaining hole
reserve is at least \(h'\sqrt L/20\), which absorbs \(c_*h'\) for large
\(R\). Also
\[
 (s_1+s_2)p^2\log(s+1)
 \le2s^{4/5+2a}\log(s+1)=2s^{0.82}\log(s+1).
\]
For fixed \(q\), the corresponding \(O_q\) error is at most \(s^{9/10}\)
once \(R\) is large.
These comparisons prove \eqref{ac:eq:21} for every allowed
\(\mathcal H'\) and hook type, including the types of small dimension.

\subsection{Adding the cells outside the hook}
\label{ac:part:245}

Return to the diagram \(\lambda\) in \eqref{ac:eq:17}. Keep its
\((p,p)\)-hook part \(\gamma\), which is a Young diagram. Translate the
remaining boxes by subtracting \(p\) from both indices to obtain an
ordinary partition \(\delta\), and put \(k=|\delta|\). The skew shape
\(\lambda/\gamma\) is a translate of \(\delta\), so branching and
Frobenius reciprocity give multiplicity \(D_\delta\) for \(\lambda\) in the representation
induced from \([\gamma]\) on \(S_{M-k}\) to \(S_M\). That induced space has
dimension \((M)_kD_\gamma\). In particular,
\begin{equation}
 F(\lambda)\le F(\gamma)+k\log s.
 \label{ac:eq:22}
\end{equation}

Realize the induced space with one fiber of type \([\gamma]\) for each
placement of \(k\) distinguishable cards in the free slots. A bijection
moves the placement and acts on its fiber by the bijection of the
remaining slots; consistent identifications of those slot sets change
the fiber maps only by unitaries. The average block from placement \(x\)
to placement \(y\) is \(p_{\mathcal H}(x,y)\) times the conditional
operator in \([\gamma]\) after those extra trajectories have been
prescribed. If this probability is positive, the endpoints force a
feasible augmented family \(\mathcal H'\) of \(h+k\) trajectories.
Thus \eqref{ac:eq:21} applies to every nonzero block, with no dimension
condition on \(\gamma\).

The \(2q\)-moment of the induced average is the sum of the irreducible
moments with their multiplicities, and hence is at least
\(D_\delta T_\lambda^{\mathcal H}\). Expand its trace as a sum over
cyclic placement indices with \(2q\) edges, alternating between the
input and output placement sets. Let \(p_j\) be the transition entry on
edge \(j\), using the transposed transition matrix on an adjoint edge.
The fiber map on an adjoint edge is the adjoint of its corresponding
forward conditional operator. It has the same singular values and uses
the same forward path potential. For a nonzero cyclic term, write
\(\mathcal H'_j\) for these augmented paths on edge \(j\).
The triangle inequality and matrix trace H\"older give
\begin{equation}
 D_\delta T_\lambda^{\mathcal H}
 \le \sum
 \left(\prod_{j=1}^{2q}p_j\right)
 \prod_{j=1}^{2q}(T_\gamma^{\mathcal H'_j})^{1/(2q)}.
 \label{ac:eq:23}
\end{equation}
The sum may be restricted to nonzero cyclic terms.

Apply \eqref{ac:eq:21}. Apart from the potential
\(-C(\mathcal H'_j)\), its terms depend only on \(s,h+k,\gamma\), so the
same bound applies on every edge. To handle the potentials, apply
\eqref{ac:eq:4} to a \(1/(2q)\) power of each \(p_j\). Multiplication
over the edges yields
\[
 \prod_{j=1}^{2q}\left(p_j e^{-C(\mathcal H'_j)/(2q)}\right)
 \le s^{-k}\exp\{-C(\mathcal H)+(\log4)kb\}
       \prod_{j=1}^{2q}p_j^{1-1/(2q)}.
\]
The augmented potentials have disappeared, leaving the potential of the
original path family.

It remains to sum the last product. Put \(n_k=(M)_k\) and label the
\(2q\) placement variables \(x_1,\ldots,x_{2q}\) cyclically. Each edge
matrix is the placement transition matrix or its transpose. Both row
and column sums are one, because the transition averages bijections of
the free slots. For each \(r\), let \(A_r\) be the product of all edge
entries except the \(r\)th. Then
\[
 \prod_{j=1}^{2q}p_j^{1-1/(2q)}
 =\prod_{r=1}^{2q}A_r^{1/(2q)}.
\]
Deleting an edge opens the cycle into a chain. Successive summation over
an endpoint of that chain uses a row or column sum equal to one; the
last placement has \(n_k\) possible values. Hence \(\sum A_r=n_k\)
for every \(r\), and scalar H\"older gives
\[
 \sum_{x_1,\ldots,x_{2q}}\prod_{j=1}^{2q}p_j^{1-1/(2q)}
 \le\prod_{r=1}^{2q}\left(\sum A_r\right)^{1/(2q)}
 =n_k.
\]
Restricting to the nonzero cyclic terms only decreases this sum.
The net counting factor in \eqref{ac:eq:23} is therefore
\(s^{-k}(M)_k\le1\). We have proved
\begin{equation}
 \log T_\lambda^{\mathcal H}
 \le-\widehat c F(\gamma)+\widehat e(h+k)\log s+s^{9/10}
       -C(\mathcal H)+(\log4)kb-F(\delta).
 \label{ac:eq:24}
\end{equation}

We finish by comparing the cost of the extra trajectories with the
dimension gained from the removed boxes and from the hook estimate.
Since \(\widehat c>0\), \eqref{ac:eq:22} gives
\[
\begin{aligned}
 \log T_\lambda^{\mathcal H}
 \le{}&-\widehat c F(\lambda)+\widehat e h\log s
       -C(\mathcal H)+s^{9/10}\\
 &+\big[(\widehat c+\widehat e)k\log s
       +(\log4)kb-F(\delta)\big].
\end{aligned}
\]
The positive \(k\)-terms in the bracket are the cost of changing the
dimension term from \(\gamma\) to \(\lambda\) and prescribing \(k\)
extra trajectories. The coefficient of \(k\log s\) is
\[
 \widehat c+\widehat e=c_0+e_0=\frac a{50}.
\]

Every row of \(\delta\) is no longer than the \((p+1)\)st row of
\(\lambda\), and likewise for columns. Those two lengths are at most
\(s/(p+1)\) by the total size of \(\lambda\). Thus every hook of \(\delta\) has length at
most \(2s/(p+1)\). If \(k\ge s^{1-a/2}\), the hook formula and
\(k!\ge(k/e)^k\) give
\[
 F(\delta)\ge
 k\log\frac{k(p+1)}{2es}
 \ge k\left(\frac a2\log s-\log(2e)\right).
\]
This dominates the bracket's positive \(k\)-terms for large \(R\),
because \(a/50<a/2\) and \(b/\log s\le1/\log R\). If instead
\(k<s^{1-a/2}\), discard \(-F(\delta)\le0\); the bracket is then at
most \(O(s^{1-a/2}\log s)\).

Finally, \eqref{ac:eq:17} gives
\[
 (\widehat c-c(s))F(\lambda)
 \ge\frac{s^{1-a/4}}{20\sqrt{\log s}},
 \qquad
 s^{9/10}+s^{1-a/2}\log s
 =o\!\left(\frac{s^{1-a/4}}{\sqrt{\log s}}\right).
\]
For sufficiently large \(R\), the dimension reserve therefore absorbs
the remaining errors. Since \(\widehat e<e(s)\), the old hole allowance
also improves. This proves \eqref{ac:eq:15}.

The induction is complete. All scale requirements, including the error
absorptions for every \(s\ge R\), hold after fixing a sufficiently large
absolute power of two \(R\) after \(q\). The final small-real-\(z\)
interval is obtained by choosing \(z_*>0\) for \eqref{ac:eq:3}, the sparse
estimate, and the finite base case.
